% !TeX program = xelatex
% =============================================================
% 模板 01：中文课程论文（ctex + XeLaTeX + BibTeX / GB/T 7714）
% 编译： scripts\build.ps1 cn-course-paper
% 需要改： 标题、作者、单位、摘要、关键词、正文、refs.bib
% 版式： A4，上下 2.54cm，左右 3.17cm，正文小四，1.5 倍行距
% =============================================================
\documentclass[UTF8,a4paper,zihao=-4,fontset=windows]{ctexart}

\usepackage{geometry}
\geometry{top=2.54cm,bottom=2.54cm,left=3.17cm,right=3.17cm}
\usepackage{amsmath,amssymb,bm}
\usepackage{booktabs}
\usepackage{array}
\usepackage{multirow}
\usepackage{graphicx}
\usepackage{caption}
\usepackage{gbt7714}
\usepackage[hidelinks]{hyperref}

\linespread{1.5}
\pagestyle{plain}   % 页码居中页脚，不显示页眉
\captionsetup{font=small,labelsep=quad}
\captionsetup[table]{position=top}
\captionsetup[figure]{position=bottom}

% 需要代码块 / 算法伪代码 / 子图 / 跨页长表 / 图表按节编号时：
% 见 references\latex-blocks.md，复制对应导言区与正文块即可

\begin{document}

% ------------------------- 标题区（改这里） -------------------------
\begin{center}
  {\zihao{3}\heiti 占位标题：基于 XXX 方法的 YYY 问题研究}\\[0.8em]
  {\zihao{-4} 张\quad 三}\\[0.4em]
  {\zihao{5}（XX 大学 \ XX 学院，XX 专业，学号 2026000000）}
\end{center}

\vspace{0.5em}

% ------------------------- 摘要（改这里） -------------------------
\noindent{\heiti 摘\quad 要：}这里是摘要占位文字。本文围绕 XXX 问题，采用 YYY 方法建立了 ZZZ 模型，并通过实验验证了模型的有效性。摘要应独立成篇，200--300 字，不出现图表和公式编号。

\vspace{0.3em}
\noindent{\heiti 关键词：}占位关键词一；占位关键词二；占位关键词三

\vspace{0.6em}
\hrule
\vspace{0.8em}

% ------------------------- 正文 -------------------------
\section{引言}

这里是正文占位内容。课程论文的引言通常交代研究背景、国内外现状和本文工作。引用参考文献时使用 \verb|\cite|，例如文献~\cite{zhang2024demo} 给出了相关结论，英文文献示例见~\cite{smith2023demo}。

\section{方法}

\subsection{模型描述}

设输入变量为 $x_i$，权重为 $\alpha_i$，则加权综合指标可写为
\begin{equation}
  S(x)=\sum_{i=1}^{n}\alpha_i x_i+\beta_0,
  \label{eq:weighted}
\end{equation}
其中 $\beta_0$ 为偏置项。式~\eqref{eq:weighted} 为公式占位示例，可直接替换。

\subsection{参数设置}

实验参数见表~\ref{tab:params}。

% ------------------------- 表格 -------------------------
\begin{table}[htbp]
  \centering
  \caption{实验参数设置（表格占位示例）}
  \label{tab:params}
  \begin{tabular}{lccc}
    \toprule
    参数 & 符号 & 取值 & 单位 \\
    \midrule
    学习率     & $\eta$   & 0.01 & --   \\
    迭代次数   & $N$      & 1000 & 次   \\
    正则系数   & $\lambda$& 0.10 & --   \\
    \bottomrule
  \end{tabular}
\end{table}

\section{结果与分析}

表~\ref{tab:params} 给出参数配置，按式~\eqref{eq:weighted} 计算后可对比不同方案的效果。此处补充结果描述、图表解释与误差分析。

\section{结论}

本文完成了占位内容的研究，得到如下结论：（1）\ldots；（2）\ldots。

% ------------------------- 参考文献 -------------------------
\bibliographystyle{gbt7714-numerical}
\bibliography{refs}

\end{document}
